\chapter{2014年上海卷（春）}

\section{填空题}

\begin{problem}
若 \( 4^{x}=16 \)，则 \(x=\)\fillinblank{}.

\end{problem}

\begin{problem}
计算：\( \i(1+\i)= \)\fillinblank{}（i 为虚数单位）.

\end{problem}

\begin{problem}
1, 1, 2, 2, 5 这五个数的中位数是 \fillinblank{}.

\end{problem}

\begin{problem}
若函数 \( f(x)=x^{3}+a \) 为奇函数，则实数 \( a=\)\fillinblank{}.

\end{problem}

\begin{problem}
点 \( O(0,0) \) 到直线 \( x + y - 4 = 0 \) 的距离是 \fillinblank{}.

\end{problem}

\begin{problem}
函数 \( y=\frac{1}{x+1} \) 的反函数为 \fillinblank{}.

\end{problem}

\begin{problem}
已知等差数列 \( \left\{a_{n}\right\} \) 的首项为 1，公差为 2，则该数列的前 \(n\) 项和 \( S_{n}= \)\fillinblank{}.

\end{problem}

\begin{problem}
已知 \( \cos\alpha=\frac{1}{3} \)，则 \( \cos2\alpha= \)\fillinblank{}.

\end{problem}

\begin{problem}
已知 \( a \)、\( b \in R^n \). 若 \( a + b = 1 \)，则 \( ab \) 的最大值是 \fillinblank{}.

\end{problem}

\begin{problem}
在 10 件产品中，有 3 件次品，从中随机取出 5 件，则恰含 1 件次品的概率是 \fillinblank{}（结果用数值表示）.

\end{problem}

\begin{problem}
某货船在 O 处看灯塔 M 在北偏东 \( 30^{\circ} \) 方向，它以每小时 18 海里的速度向正北方向航行，经过 40 分钟到达 B 处，看到灯塔 M 在北偏东 \( 75^{\circ} \) 方向，此时货船到灯塔 M 的距离为 \fillinblank{}~海里.

\end{problem}

\begin{problem}
已知函数 \( f(x)=\frac{x-2}{x-1} \) 与 \( g(x)=mx+1-m \) 的图像相交于 A、B 两点. 若动点 P 满足 \( \left|\overrightarrow{PA}+\overrightarrow{PB}\right|=2 \)，则 P 的轨迹方程为 \fillinblank{}.

\end{problem}

\section{单选题}

\begin{problem}
两条异面直线所成的角的范围是（\quad）

\choices
    {\( \left(0, \frac{\pi}{2}\right) \)}
    {\( \left(0, \frac{\pi}{2}\right] \)}
    {\( \left[0, \frac{\pi}{2}\right) \)}
    {\( \left[0, \frac{\pi}{2}\right] \)}

\end{problem}

\begin{problem}
复数 \( 2+\i \)（i 为虚数单位）的共轭复数为（\quad）

\choices
    {2 - i}
    {\(-2\) + i}
    {\(-2\) - i}
    {\( 1 + 2\i \)}

\end{problem}

\begin{problem}
右图是下列函数中某个函数的部分图像，则该函数是（\quad）

\choices
    {\( y = \sin x \)}
    {\( y = \sin 2x \)}
    {\( y = \cos x \)}
    {\( y = \cos 2x \)}

\end{problem}

\begin{problem}
在 \( (x+1)^{4} \) 的二项展开式中，\( x^{2} \) 项的系数为（\quad）

\choices
    {6}
    {4}
    {2}
    {1}

\end{problem}

\begin{problem}
下列函数中，在 \( \R \) 上为增函数的是（\quad）

\choices
    {\( y = x^{2} \)}
    {\( y = |x| \)}
    {\( y = \sin x \)}
    {\( y = x^{3} \)}

\end{problem}

\begin{problem}
\( \left|\begin{array}{cc}\cos\theta & -\sin\theta \\ \sin\theta & \cos\theta\end{array}\right| = \)（\quad）

\choices
    {\( \cos 2\theta \)}
    {\( \sin 2\theta \)}
    {1}
    {\(-1\)}

\end{problem}

\begin{problem}
设 \( x_{0} \) 为函数 \( f(x)=2^{x}+x-2 \) 的零点，则 \( x_{0}\in \)（\quad）

\choices
    {\( (-2,-1) \)}
    {\( (-1,0) \)}
    {\( (0,1) \)}
    {\( (1,2) \)}

\end{problem}

\begin{problem}
若 \( a > b \)，\( c \in \R \)，则下列不等式中恒成立的是（\quad）

\choices
    {\( \frac{1}{a}<\frac{1}{b} \)}
    {\( a^{2}>b^{2} \)}
    {\( a|c|>b|c| \)}
    {\( \frac{a}{c^{2}+1}>\frac{b}{c^{2}+1} \)}

\end{problem}

\begin{problem}
若两个球的体积之比为 8:27，则它们的表面积之比为（\quad）

\choices
    {2:3}
    {4:9}
    {8:27}
    {\( 2\sqrt{2}:3\sqrt{3} \)}

\end{problem}

\begin{problem}
已知数列 \( \left\{a_{n}\right\} \) 是以 \(q\) 为公比的等比数列，若 \( b_{n}=-2a_{n} \)，则数列 \( \left\{b_{n}\right\} \) 是（\quad）

\choices
    {以 \(q\) 为公比的等比数列}
    {以 \(-q\) 为公比的等比数列}
    {以 \(2q\) 为公比的等比数列}
    {以 \(-2q\) 为公比的等比数列}

\end{problem}

\begin{problem}
若点 P 的坐标为 \( (a, b) \)，曲线 C 的方程为 \( F(x, y) = 0 \)，则“ \( F(a, b) = 0 \)”是“点 P 在曲线 C 上”的（\quad）

\choices
    {充分非必要条件}
    {必要非充分条件}
    {充分必要条件}
    {既非充分又非必要条件}

\end{problem}

\begin{problem}
如图，在底面半径和高均为 1 的圆锥中，AB、CD 是底面圆 O 的两条互相垂直的直径，E 是母线 PB 的中点. 已知过 CD 与 E 的平面与圆锥侧面的交线是以 E 为顶点的抛物线的一部分，则该抛物线的焦点到圆锥顶点 P 的距离为（\quad）

\choices
    {1}
    {\( \frac{\sqrt{3}}{2} \)}
    {\( \frac{\sqrt{6}}{2} \)}
    {\( \frac{\sqrt{10}}{4} \)}

\end{problem}

\section{解答题}

\begin{problem}
已知不等式 \( \frac{x-2}{x+1}<0 \) 的解集为 \(A\)，函数 \( y=\lg(x-1) \) 的定义域为集合 \(B\)，求 \( A\cap B \).

\end{problem}

\begin{problem}
已知函数 \( f(x)=x^{2}-4x+a, x \in [-3,3] \). 若 \( f(1)=2 \)，求 \( y=f(x) \) 的最大值和最小值.

\end{problem}

\begin{problem}
如图，在体积为 \( \frac{1}{3} \) 的三棱锥 \(P-ABC\) 中，PA 与平面 ABC 垂直，AP = AB = 1，\( \angle BAC=\frac{\pi}{2} \)，E、F 分别是 PB、AB 的中点. 求异面直线 EF 与 PC 所成的角的大小（结果用反三角函数值表示）.

\end{problem}

\begin{problem}
已知椭圆 \( C: \frac{x^{2}}{a^{2}} + y^{2} = 1 (a > 1) \) 的左焦点为 F，上顶点为 B. （1）若直线 FB 的一个方向向量为 \( (1, \frac{\sqrt{3}}{3}) \)，求实数 \(a\) 的值；（2）若 \( a = \sqrt{2} \)，直线 \( l: y = kx - 2 \) 与椭圆 \( C \) 相交于 \( M \)、\( N \) 两点，且 \( \overrightarrow{FM} \cdot \overrightarrow{FN} = 3 \)，求实数 \(k\) 的值.

\end{problem}

\begin{problem}
已知数列 \( \left\{a_{n}\right\} \) 满足 \( a_{n}>0 \)，双曲线 \( C_{n}:\frac{x^{2}}{a_{n}}-\frac{y^{2}}{a_{n+1}}=1(n\in N^{*}) \). （1）若 \( a_{1}=1, a_{2}=2 \)，双曲线 \( C_{n} \) 的焦距为 \( 2c_{n} \)，\( c_{n}=\sqrt{4n-1} \)，求 \( \left\{a_{n}\right\} \) 的通项公式；（2）如图，在双曲线 \( C_n \) 的右支上取点 \( P_n(x_n, n) \)，过 \( P_n \) 作 \( y \) 轴的垂线，在第一象限内交 \( C_n \) 的渐近线于点 \( Q_n \)，联结 \( OP_n \)，记 \( \Delta OP_nQ_n \) 的面积为 \( S_n \). 若 \( \lim_{n \to \infty} a_n = 2 \)，求 \( \lim_{n \to \infty} S_n \). （关于数列极限的运算，还可参考如下性质：若 \( \lim_{n\to\infty}u_n=A(u_n\geq0) \)，则 \( \lim_{n\to\infty}\sqrt{u_n}=\sqrt{A} \)）

\end{problem}

\begin{problem}
已知直角三角形 ABC 的两直角边 AC、BC 的边长分别为 \(b, a\)，如图，过 AC 边的 \(n\) 等分点 \( A_{i} \) 作 AC 边的垂线 \( d_{i} \)，过 BC 边的 \(n\) 等分点 \( B_{i} \) 和顶点 A 作直线 \( l_{i} \)，记 \( d_{i} \) 与 \( l_{i} \) 的交点为 \( P_{i}(i=1,2,\cdots,n-1) \). 是否存在一条圆锥曲线，对任意的正整数 \( n \geq 2 \)，点 \( P_{i}(i = 1, 2, \cdots, n-1) \)

\end{problem}

\begin{problem}
某人造卫星在地球赤道平面绕地球飞行，甲、乙两个监测点分别位于赤道上东经 \( 131^{\circ} \) 和 \( 147^{\circ} \)，在某时刻测得甲监测点到卫星的距离为 1537.45 千米，乙监测点到卫星的距离为 887.64 千米. 假设地球赤道是一个半径为 6378 千米的圆，求此时卫星所在位置的高度（结果精确到 0.01 千米）和经度（结果精确到 0.01°）.

\end{problem}

\begin{problem}
如果存在非零常数 \(c\)，对于函数 \( y = f(x) \) 定义域 \(R\) 上的任意 \(x\)，都有 \( f(x + c) > f(x) \) 成立，那么称函数为“Z 函数”. （1）求证：若 \( y = f(x) (x \in R) \) 是单调函数，则它是“Z 函数”；（2）若函数 \( g(x)=ax^{3}+bx^{2} \) 是“Z 函数”，求实数 \(a,b\) 满足的条件.
\end{problem}

